%! AP Statistics — Unit 3, Lesson 3.3 — Homework KEY
\documentclass[10pt]{article}
\usepackage{apstats-article}
\usepackage{apstats-key}

\begin{document}

\pageheader{Unit 3, Lesson 3.3}{Homework — KEY}
\namedateperiod

\vspace{0.1in}

\begin{practicebox}
\small
\textbf{1.}
\begin{enumerate}[leftmargin=1.5em, itemsep=5pt]
  \item \ans{Simple random sample (SRS)}
  \item \ans{Cluster sample}
  \item \ans{Systematic random sample}
  \item \ans{Stratified random sample}
\end{enumerate}
\end{practicebox}

\vspace{0.08in}

\begin{practicebox}
\small
\textbf{2.}
\begin{enumerate}[leftmargin=1.5em, itemsep=5pt]
  \item \ans{Divide students into 5 strata (one per school). Randomly select 20 students from each school using an SRS.}
  \item \ans{The 5 schools are the clusters. Randomly select 1 school (or as needed to get about 100 students); survey all students in the selected school(s).}
  \item \ans{Clusters — each school likely contains students of all types (different backgrounds, ability levels, etc.), making each school heterogeneous internally. Strata would be homogeneous groups.}
\end{enumerate}
\end{practicebox}

\vspace{0.08in}

\begin{practicebox}
\small
\textbf{3.} \ans{Sample response: Strata could be grade levels (9th, 10th, 11th, 12th). Stratifying is appropriate because physical activity levels may differ across grade levels (e.g., 9th graders may have mandatory PE; 12th graders may not). Procedure: Obtain the enrollment list for each grade. Use a random number generator to select students proportionally from each grade level until the desired sample size is reached.}
\end{practicebox}

\end{document}
